% !TEX root = ../main.tex
% ---------------------------------------------------------------
% GENERATED FILE - DO NOT EDIT.
% Produced by docs/tools/render_reference.py from the repository source.
% Regenerate with: make docs
% ---------------------------------------------------------------

\apibreak
\section{\texttt{extract\_entities}}
\label{func:feedback_intelligence_agent.memory.extract_entities}\index{Functions!extract\_entities}

Extract salient content tokens from text, preserving first-seen order.
\apifield{Usage}
\begin{lstlisting}[style=usage, language=Python]
extract_entities(text: str, *, limit: int = 4) -> list[str]
\end{lstlisting}
\apifield{Arguments}
\begin{arglist}
  \item[\texttt{text}] \texttt{str}, required.
  \item[\texttt{limit}] \texttt{int}, default \texttt{4}, keyword-only.
\end{arglist}
\apifield{Value}\texttt{list[str]\allowbreak{}}
\apifield{Details}Tokens are lowercased, at least three characters long, and filtered against a small stopword list, so the result approximates the entities and topics mentioned in the previous turn.
\apifield{Source}\texttt{src/\allowbreak{}feedback\_\allowbreak{}intelligence\_\allowbreak{}agent/\allowbreak{}memory.\allowbreak{}py:241}~(\href{https://github.com/DiogoRibeiro7/feedback-intelligence-agent/blob/edc9d78403eeed9f7d6ee5aef6b8d9b1b5780ef4/src/feedback_intelligence_agent/memory.py#L241}{online}), module \cref{module:feedback_intelligence_agent.memory}
